\chapter{Fixed-target production}
\label{ch:production}

\chaptergoals{6 of 11 in Part~II: the first simulation.}
{\cref{ch:geometry} (\file{ship\_geometry.db} in \file{work/}) and
\cref{ch:datamodel}.}
{generate proton--tungsten collisions with \pkg{aegir}, record what leaves the
hadron absorber on a scoring plane, normalise counts to protons on target, and
look inside the file you made.}

\section{What aegir is}

\repo{aegir} is the SHiP simulation: a set of Phlex plugins that generate or
read particles, transport them through the detector with Geant4, and write
the result. A simulation job is a graph of four kinds of node:

\begin{center}
\begin{tikzpicture}[
  every node/.style={font=\sffamily\small},
  box/.style={draw=shipnavy,fill=codebg,rounded corners=2pt,
    text width=35mm,minimum height=12mm,align=center},
  job/.style={box,draw=shipblue,fill=shipblue!6,text width=55mm},
  flow/.style={-{Stealth[length=2mm]},line width=0.7pt},
  uses/.style={flow,dashed,color=shipblue}]
  \node[box] (src) at (0,0) {source\\\scriptsize\texttt{fixed\_target\_source}};
  \node[box] (g4) at (5.2,0) {transform\\\scriptsize\texttt{geant4\_module}};
  \node[box] (out) at (10.4,0) {observer\\\scriptsize\texttt{sim\_output\_module}};
  \node[job] (geo) at (2.1,2.4) {job provider\\\scriptsize\texttt{geometry\_geomodel\_provider}};
  \node[job] (fld) at (8.4,2.4) {job provider\\\scriptsize\texttt{field\_null\_provider}};
  \draw[flow] (src) -- node[above,font=\scriptsize]{particles} (g4);
  \draw[flow] (g4) -- node[above,font=\scriptsize]{sim\_result} (out);
  \draw[uses] (geo.south) -- ([xshift=-7mm]g4.north);
  \draw[uses] (fld.south) -- ([xshift=7mm]g4.north);
  \draw[flow] (src.south) -- ++(0,-0.8) -|
    node[pos=0.25,below,font=\scriptsize]{generator record + event header} (out.south);
\end{tikzpicture}
\end{center}

Sources work per event; the geometry and the field are providers at the job
level, built once and shared by all events. The plugins available are:

\begin{center}
\small
\begin{tabularx}{\linewidth}{l l L}
\toprule
Plugin & Kind & Purpose \\
\midrule
\texttt{particle\_gun\_source}       & source   & single particles in a cone \\
\texttt{pythia8\_source}             & source   & Pythia\,8 fixed-target p--p collisions \\
\texttt{fixed\_target\_source}       & source   & \textbf{Pythia\,8 on tungsten, with the interaction point sampled in the target (stage 1)} \\
\texttt{genie\_reader\_source}       & source   & pre-generated GENIE neutrino interactions \\
\texttt{file\_source}                & source   & \textbf{particles read from a data-model file (stage 2)} \\
\texttt{geometry\_builtin\_provider} & provider & a test geometry (tungsten block and silicon planes) \\
\texttt{geometry\_gdml\_provider}    & provider & geometry from a GDML file \\
\texttt{geometry\_geomodel\_provider}& provider & \textbf{the SHiP geometry via the geometry service} \\
\texttt{field\_null\_provider}       & provider & no field \\
\texttt{field\_map\_provider}        & provider & field maps via the field service \\
\texttt{geant4\_module}              & transform& Geant4 transport, multi-threaded \\
\texttt{sim\_output\_module}         & observer & RNTuple output and validation histograms \\
\bottomrule
\end{tabularx}
\end{center}

\section{Protons on tungsten}

\texttt{fixed\_target\_source} generates one 400\,GeV proton interaction per
event. Tungsten has 74 protons and 110 neutrons, so it runs two Pythia\,8
instances, one for p--p and one for p--n collisions, and picks one in that
proportion. The interaction point is drawn along the target from the
exponential law with the interaction length of tungsten, 19.2\,cm, between
$z=0$ and the end of the target at 1.164\,m. Pythia decays only particles
whose proper decay length $c\tau_0$ is below a threshold (1\,mm); everything
longer-lived, pions and kaons included, is handed to Geant4, which then
decides whether it decays or interacts in the tungsten.

Only the status-1 entries of the generator record are tracked
(\cref{sec:mcstatus}); with \texttt{record: 'final\_state'} these are the only
entries written.

\section{The scoring plane}
\label{sec:scoringplane}

Simulating every proton through the whole of SHiP would be wasteful: nearly
all particles stop in the target or in the hadron absorber, 2.3\,m of iron
behind it. Stage~1 therefore stops at a \emph{scoring plane} 1\,cm behind the
absorber, at $z\approx4.36$\,m, in the gap before the first muon-shield
magnet. The plane is the one FairShip calls \texttt{exitHadronAbsorber}:
7.12\,m wide and 3.40\,m high, wider than the muon shield itself, so that it
also catches what passes beside it.

\begin{description}[style=nextline]
  \item[It lives in a parallel world.] The plane is not part of the SHiP
    geometry; Geant4 overlays it, so it can cut through any volume and
    nothing in \file{ship\_geometry.db} has to change.
  \item[It records everything.] Every track that reaches it, charged or
    neutral, deposits energy or not, becomes one \texttt{SimHit} with
    \field{detector\_id}~100 and the plane's id in \field{geometry\_node\_id}.
  \item[It kills.] With \texttt{kill: true}, a track is stopped at the plane,
    so the job ends there.
  \item[It hands over to stage 2.] With \texttt{record\_particles: true}, the
    tracks it records going downstream are also written as
    \texttt{plane\_particles}: as they were when they entered the plane, with
    their ancestry. SHiP's stage~2 simulates them onwards from there; you
    can try it in \cref{ex:ownreplay}.
\end{description}

Tracks with less kinetic energy than the cut \texttt{ecut} (0.5\,GeV, as in
FairShip) are killed wherever they are: they cannot reach the decay volume
anyway. Of the tracked particles only those that reach the plane and their
ancestors are written (\texttt{keep\_particles: 'hit\_ancestry'}), and events
in which nothing reached the plane are not written at all
(\texttt{filter\_empty}).

\section{The workflow}

The workflow for this stage is \file{code/workflows/ft\_production.jsonnet}.
It is the job of \pkg{aegir}'s \file{workflows/fixed\_target\_production.jsonnet}
with every building block written out, so that you can read the whole job in
one place:

\codefile[linerange={20-111},basicstyle=\ttfamily\footnotesize]{jsonnet}{code/workflows/ft_production.jsonnet}

\begin{exercise}{Read before you run}
From the workflow alone: which detector volumes are sensitive in this stage?
Which tracks are killed, and where? Which file will hold the result, and
which seed is used if you do not give one?
\end{exercise}

\begin{solution}
None: \texttt{sensitive\_volumes} is empty, and only the scoring plane records
hits. Tracks below 0.5\,GeV kinetic energy are killed anywhere; every track
reaching the plane is killed there. \file{ft\_production.root}, with the
validation histograms in \file{ft\_production\_validation.root}; seed~42.
\end{solution}

\section{Running it}

\begin{exercise}{Your first protons on target}
\label{ex:production}
Open a shell in the \pkg{aegir} environment and go to your working
directory, which must contain \file{ship\_geometry.db} (\cref{ch:geometry}):
\begin{shell}
cd "$SHIP_TUTORIAL/aegir"
pixi cvmfs
cd "$SHIP_TUTORIAL/work"
POT=1000                         # protons on target for this tutorial
time phlex -j 4 -c <(jsonnet --tla-code events=$POT --tla-code seed=42 \
    $CODE/workflows/ft_production.jsonnet)
\end{shell}
Watch the log: the geometry is read and converted, the plane is placed behind
\texttt{/SHiP/muon\_shield/magn\_absorb}, the seeds are printed, and the event
loop reports progress every 100 events. Keep the value of \texttt{POT}: every
normalisation from now on needs it.
\end{exercise}

\begin{organisernote}
In the dry run on lxplus (21 October 2026), 1000 protons took 256\,s with
\cmd{-j 4}, about 100\,s of it start-up. That many
protons give only a handful of muons at the plane, which is why
\cref{ch:aegir} continues with an artificial muon sample.
\end{organisernote}

\begin{checkpoint}
Two new files, \file{ft\_production.root} and
\file{ft\_production\_validation.root}, and no error. The file holds
\emph{fewer} entries than \texttt{POT}: only the events in which something
reached the plane.
\end{checkpoint}

\begin{exercise}{Look inside the file}
\label{ex:openfile}
Still in the \pkg{aegir} environment, start Python:
\begin{py}
import ROOT
reader = ROOT.RNTupleReader.Open("events", "ft_production.root")
reader.PrintInfo()     # fields, types, number of entries, compression
reader.Show(0)         # the first entry, as JSON
\end{py}
Which fields does the file have (compare with \cref{sec:filelayout})? In the
first entry, how many generator particles are there, and how many hits? What
is the \field{detector\_id} of the hits?
\end{exercise}

\begin{solution}
The six fields of the simulation output: \texttt{event\_number},
\texttt{event\_header}, \texttt{mc\_particles}, \texttt{sim\_hits},
\texttt{sim\_particles}, \texttt{plane\_particles}. A p--W interaction gives
some tens of final-state particles in \texttt{mc\_particles}. The hits are
all plane hits, \field{detector\_id}~100 with \field{geometry\_node\_id}~1,
one per track reaching the plane, and there are as many status-1
\texttt{plane\_particles}, plus their ancestors with other statuses.
\end{solution}

\section{What leaves the absorber}

\file{code/analysis/plane\_summary.py} reads the plane hits, counts them by
particle type, and normalises the counts to the protons you simulated, per
proton and per spill of $4\times10^{13}$ protons:

\begin{exercise}{Count what reaches the plane}
\label{ex:planesummary}
\begin{shell}
python $CODE/analysis/plane_summary.py ft_production.root --pot $POT
\end{shell}
It prints a table and writes three PDFs: the muon momentum and transverse
momentum, and where the muons cross the plane.
\begin{enumerate}
  \item Which particles reach the plane, and which are most numerous?
  \item How many muons per proton on target, and per spill? Are there more
        $\mu^+$ or $\mu^-$, and why?
  \item Look at \file{plane\_muon\_p.pdf}: where are most muons, and how far
        does the spectrum reach?
  \item Why does the file need \cmd{--pot}, rather than counting its entries?
\end{enumerate}
\end{exercise}

\begin{solution}
\begin{enumerate}
  \item Neutrinos dominate, then neutrons and muons: the pion and kaon
        decays that make the muons make the neutrinos too, and nothing stops
        the neutrinos. In the dry run (seed 42), 31 of the 1000 protons left
        something at the plane: 23 neutrinos, 5 neutrons and 3 muons.
  \item In the dry run, 3 muons in 1000 protons, $3\times10^{-3}$ per proton
        or about $10^{11}$ per spill, the flux the muon shield has to deal
        with. With three muons the uncertainty is about 60\,\%. Expect more
        $\mu^+$ than $\mu^-$ in a large sample: the beam and the target carry
        positive charge, so positive pions and kaons, and their decay muons,
        are more numerous than negative ones. The dry run had one $\mu^+$ and
        two $\mu^-$, which says nothing either way.
  \item Most muons have a few GeV, just above the cut; the spectrum falls
        steeply and reaches some hundred GeV. With a small sample the tail is
        empty: that is a statistics effect, not a cut.
  \item The output holds only events with something at the plane
        (\texttt{filter\_empty}); the protons whose products all stopped are
        not in the file, but they are part of the normalisation.
\end{enumerate}
\end{solution}

\begin{exercise}{The validation histograms}
\pkg{aegir} fills histograms as it runs. Open them with a ROOT browser:
\begin{shell}
root -l ft_production_validation.root
root [1] new TBrowser
\end{shell}
Find \texttt{h\_plane\_hit\_pdg} and \texttt{h\_plane\_hit\_momentum}, and
compare them with your table. What does \texttt{h\_mc\_multiplicity} show?
\end{exercise}

\begin{solution}
\texttt{h\_plane\_hit\_pdg} has peaks at $\pm 12$, $\pm 14$ (neutrinos) and
$\pm 13$ (muons), with the same counts as the table. \texttt{h\_mc\_multiplicity}
is the number of final-state generator particles per event, for the written
events only.
\end{solution}

\begin{exercise}[colback=white]{Optional: the energy cut}
Run again with \cmd{--tla-code ecut=1.0} and \cmd{--tla-str output=ft\_ecut1},
and compare the muon tables. Which muons went missing, and does it matter for
the background in the decay volume?
\end{exercise}

\begin{solution}
The muons between 0.5 and 1\,GeV. They range out in the muon-shield iron long
before the decay volume, so for the background they do not matter, and a
higher cut only saves CPU time. That is the reasoning behind the cut.
\end{solution}

\begin{note}
The production stage is where most CPU time goes, which is why SHiP produces
it once, in large batches, and reuses the plane particles for many replays
with different shield designs or detector settings. The plane particles of
such a production are exactly what \file{ft\_replay.jsonnet} reads
(\cref{ex:ownreplay}).
\end{note}
